% !TEX program = xelatex
% Homework template

\documentclass[cn,12pt]{homework}
% en is for English language
% cn is for Chinese language

%----- text fonts -----
% \usepackage{newtxtext}
% \setmainfont{Times New Roman}

%----- math font -----
\usepackage{newtxmath}
\usepackage{mathptmx}
\usepackage{mathpazo}
\usepackage{listings}
\usepackage{enumitem}
\usepackage{float}
\usepackage{xcolor}
\usepackage{fancyhdr}
\usepackage{algorithm}
\usepackage{algorithmic}
\usepackage{amsmath}

\pagestyle{fancy}
\lstset{
    %backgroundcolor=\color{red!50!green!50!blue!50},%代码块背景色为浅灰色
    rulesepcolor= \color{gray}, %代码块边框颜色
    breaklines=true,  %代码过长则换行
    numbers=left, %行号在左侧显示
    numberstyle= \small,%行号字体
    keywordstyle= \color{blue},%关键字颜色
    commentstyle=\color{gray}, %注释颜色
    frame=shadowbox%用方框框住代码块
    }

%----- custom theorem -----
\newtheorem{innercustomgeneric}{\customgenericname}
\providecommand{\customgenericname}{}
\newcommand{\newcustomtheorem}[2]{%
  \newenvironment{#1}[1]
  {%
   \renewcommand\customgenericname{#2}%
   \renewcommand\theinnercustomgeneric{##1}%
   \innercustomgeneric
  }
  {\endinnercustomgeneric}
}

\newcustomtheorem{ntheorem}{定理}
\newcustomtheorem{nlemma}{引理}

%----- list style -----
\setlist{nolistsep}

% differential operator
\newcommand{\dif}{\mathop{}\!\mathrm{d}}

% new command
\newcommand{\CC}{\ensuremath{\mathbb{C}}}
\newcommand{\RR}{\ensuremath{\mathbb{R}}}
\newcommand{\A}{\mathcal{A}}
\newcommand{\bA}{\boldsymbol{A}}
\newcommand{\ii}{\mathrm{i}\,}
\newcommand{\dx}[1][x]{\mathop{}\!\mathrm{d}#1}
\newcommand{\abs}[1]{\lvert#1\rvert}
\newcommand{\norm}[1]{\left\lVert#1\right\rVert}
\newcommand{\red}[1]{\textcolor{red}{#1}}

%----------------------------------------------------
%	HOMEWORK INFORMATION
%----------------------------------------------------

\lhead{\itshape Artificial Intelligence   } % 页眉
\rhead{}
\title{HOMEWORK \#3} % 作业名字

\date{Date: \today} % 日期

\institute{ZHEJIANG UNIVERSITY\quad COLLEGE OF INFORMATION SCIENCE AND ELECTRONICS ENGINEERING} % 学院或学校

\courseinfo{Artificial Intelligence } % 课程信息

\studentinfo{Name: \textit{}  \quad  \quad Student ID: \textit{}  \quad \\ Major: \textit{Electronic Science and Technology}} % 学生信息


\begin{document}

\maketitle

%----------------------------------------------------
%	作业内容
%----------------------------------------------------

%\section*{作业题目}

%------------------------------%
%------------------------------%
\begin{problem}
In our analysis of the wumpus world, 
we used the fact that each square contains a pit with probability 0.2, 
independently of the contents of the other squares. 
Suppose instead that exactly $N/5$ pits are scattered at random among 
the $N$ squares other than [1,1]. Are the variables $P_{i,j}$ and $P_{k,l}$ 
still independent? What is the joint distribution 
$\mathbf{P}(P_{1,1},\ldots,P_{4,4})$ now? Redo the calculation 
for the probabilities of pits in [1,3] and [2,2].
(conditional on known, b as in section 13.6).


\end{problem}


\begin{solution}
  \quad

\textbf{a. Independence of Variables:}
The variables \(P_{i,j}\) and \(P_{k,l}\) are not independent when exactly \(N/5\) pits are scattered among the \(N\) squares. 
The presence of a pit in one square affects the probability of pits in other squares because the total number of pits is fixed.
For example, if we know that there is a pit in square \((i,j)\), it reduces the number of pits available for the other squares, thus affecting the probability of pits in square \((k,l)\).

\textbf{b. Joint Distribution:}
In our problem, N=15, so there are 3 pits in total.
The joint distribution places equal probability on all combinations of pits in the squares, given that there are exactly 3 pits in total and equal zero probability on all combinations that do not satisfy this condition.
Thus ,the joint distribution can be expressed as:
\[
\mathbf{P}(P_{1,1}, P_{1,2}, \ldots, P_{4,4}) = 
\begin{cases} 
\frac{1}{\binom{15}{3}}, & \text{if } P_{1,1} = 0 \text{ and exactly 3 of } P_{1,2}, \ldots, P_{4,4} \text{ are 1}, \\
0, & \text{otherwise}.
\end{cases}
\]

\textbf{c. Probabilities of Pits in [1,3] and [2,2]:}

\begin{figure}[H]
\centering
\includegraphics[width=1.0\textwidth]{figures/image1.png}
\caption{cases that satisfy the condition}
\end{figure}

as shown in the figure, there are 5 cases that satisfy the conditional on known, b as in section 13.6)

1.consider Probabilities of Pits in [1,3]. There are 3 cases that satisfy the condition of having a pit in [1,3].
\begin{itemize}
  \item Case 1: $P_{1,3} = 1$, $P_{2,2} = 1$, $P_{3,1} = 1$. totally 1 assignment.
  \item Case 2: $P_{1,3} = 1$, $P_{2,2} = 1$, and $P_{3,1} = 0$. totally $\binom{10}{1}=10$ assignments.
  \item Case 3: $P_{1,3} = 1$, $P_{2,2} = 0$, and $P_{3,1} = 1$. totally $\binom{10}{1}=10$ assignments.
\end{itemize}
There are 2 cases that not satisfy the condition of having a pit in [1,3].
\begin{itemize}
  \item Case 1: $P_{1,3} = 0$, $P_{2,2} = 1$, $P_{3,1} = 1$. totally $\binom{10}{1}=10$ assignment.
  \item Case 2: $P_{1,3} = 0$, $P_{2,2} = 1$, $P_{3,1} = 0$. totally $\binom{10}{2}=45$ assignments.
\end{itemize}
Thus, the probability of $P_{1,3}$ is given by:
\[
P(P_{1,3}) = \alpha \langle 21, 55 \rangle = \langle 0.276, 0.724 \rangle
\]
2.consider Probabilities of Pits in [2,2]. Same as above, we can calculate the probability of $P_{2,2}$.
There are $1+10+10+45=66$ cases that satisfy the condition of having a pit in [2,2].
There are $10$ cases that not satisfy the condition of having a pit in [2,2].
Thus, the probability of $P_{2,2}$ is given by:
\[
P(P_{2,2}) = \alpha \langle 66, 10 \rangle = \langle 0.868, 0.132 \rangle
\]
\end{solution}
%------------------------------%
\newpage

\begin{problem}

Consider two medical tests, A and B, for a virus. Test A is 95\% effective at recognizing the virus when it is present, but has a 10\% false positive rate (indicating that the virus is present, when it is not). Test B is 90\% effective at recognizing the virus, but has a 5\% false positive rate. The two tests use independent methods of identifying the virus. The virus is carried by 1\% of all people. Say that a person is tested for the virus using only one of the tests, and that test comes back positive for carrying the virus. Which test returning positive is more indicative of someone really carrying the virus? Justify your answer mathematically.
\end{problem}

\begin{solution}
\quad

According to the problem, we have the following information:
\begin{align*}
P(A|V) &= 0.95 \quad \text{(Test A is 95\% effective at recognizing the virus)} \\
P(A|\urcorner V) &= 0.10 \quad \text{(Test A has a 10\% false positive rate)} \\
P(B|V) &= 0.90 \quad \text{(Test B is 90\% effective at recognizing the virus)} \\
P(B|\urcorner V) &= 0.05 \quad \text{(Test B has a 5\% false positive rate)} \\
P(V) &= 0.01 \quad \text{(The virus is carried by 1\% of all people)} \\
\end{align*}
According to Bayes' theorem, we can calculate the posterior probabilities of having the virus given a positive test result for each test.
\begin{align*}
P(V|A) &= \frac{P(A|V) \cdot P(V)}{P(A|V) \cdot P(V) + P(A|\urcorner V) \cdot P(\urcorner V)} \\
&= \frac{0.95 \cdot 0.01}{0.95 \cdot 0.01 + 0.10 \cdot 0.99} \\
&= \frac{0.0095}{0.0095 + 0.099} \\
&= \frac{0.0095}{0.1085} \\
&\approx 0.08756 \\
\end{align*}
\begin{align*}
P(V|B) &= \frac{P(B|V) \cdot P(V)}{P(B|V) \cdot P(V) + P(B|\urcorner V) \cdot P(\urcorner V)} \\
&= \frac{0.90 \cdot 0.01}{0.90 \cdot 0.01 + 0.05 \cdot 0.99} \\
&= \frac{0.0090}{0.0090 + 0.0495} \\
&= \frac{0.0090}{0.0585} \\
&\approx 0.1538 \\
\end{align*}
Therefore, the posterior probability of having the virus given a positive test result for Test A is approximately 0.0875, while the posterior probability of having the virus given a positive test result for Test B is approximately 0.1538. Since the posterior probability for Test B is higher, it is more indicative of someone really carrying the virus.



\end{solution}





\newpage

\begin{problem}
The Markov blanket of a variable is defined on page 517. Prove that avariable is independent of all other variables in the network, given its Markovblanket and derive Equation (14.12)(page 538).
\end{problem}

\begin{solution}
  \quad


\textbf{a. Derivation of Equation (14.12):}

 Let $ \mathbf{Y} = Y_1, \ldots, Y_\ell \text{ be the children of } X_i \text{ and let } \mathbf{Z}_j \text{ be the parents of } Y_j \text{ other than } X_i. $ 

\text{ Then we have}
 \begin{align*}
&\mathbf{P}(X_i | \text{MB}(X_i)) \\
&\quad = \mathbf{P}(X_i | \text{Parents}(X_i), \mathbf{Y}, \mathbf{Z}_1, \ldots, \mathbf{Z}_\ell) 
\end{align*}

The $\text{MB}(X_i)$ denotes the values of the variables in $X_i$’s Markov blanket

Then, we use Bayes’ rule to write

 \begin{align*}
&\mathbf{P}(X_i | \text{Parents}(X_i), \mathbf{Y}, \mathbf{Z}_1, \ldots, \mathbf{Z}_\ell) \\
&\quad = \alpha \mathbf{P}(X_i | \text{Parents}(X_i), \mathbf{Z}_1, \ldots, \mathbf{Z}_\ell) \mathbf{P}(\mathbf{Y} | \text{Parents}(X_i), X_i, \mathbf{Z}_1, \ldots, \mathbf{Z}_\ell)
\end{align*}

According to the fact that in Bayesian networks, a node is
independent of its nondescendants given its children. We can eliminate same variables in the conditioning set.
 \begin{align*}
&\alpha \mathbf{P}(X_i | \text{Parents}(X_i), \mathbf{Z}_1, \ldots, \mathbf{Z}_\ell) \mathbf{P}(\mathbf{Y} | \text{Parents}(X_i), X_i, \mathbf{Z}_1, \ldots, \mathbf{Z}_\ell) \\
&\quad = \alpha \mathbf{P}(X_i | \text{Parents}(X_i)) \mathbf{P}(\mathbf{Y} | X_i, \mathbf{Z}_1, \ldots, \mathbf{Z}_\ell) 
\end{align*}

finally, we expand  $\mathbf{P}(\mathbf{Y} | X_i, \mathbf{Z}_1, \ldots, \mathbf{Z}_\ell) $ as a product of the conditional probabilities of each child given its parents and the parents of $X_i$.
 \begin{align*}
&\alpha \mathbf{P}(X_i | \text{Parents}(X_i)) \mathbf{P}(\mathbf{Y} | X_i, \mathbf{Z}_1, \ldots, \mathbf{Z}_\ell)  \\
&\quad = \alpha \mathbf{P}(X_i | \text{Parents}(X_i)) \prod_{Y_j \in \text{Children}(X_i)} \mathbf{P}(Y_j | \text{Parents}(Y_j)) \tag{14.12} \\
\end{align*}

Which is the equation (14.12) in the book.

\textbf{b. Proof of Independence:}
To prove that a variable \(X\) is independent of all other variables in the network given its Markov blanket, we need to show that:

\begin{align*}
P(x_i | x_1, \ldots, x_{i-1}, x_{i+1}, \ldots, x_n) &= \frac{P(x_1, \ldots, x_n)}{P(x_1, \ldots, x_{i-1}, x_{i+1}, \ldots, x_n)} \\
&= \frac{P(x_1, \ldots, x_n)}{\Sigma_{x_i} P(x_1, \ldots, x_n)} \\
&= \frac{\Pi_{j=1}^n P(x_j | parents X_j)}{\Sigma_{x_i} \Pi_{j=1}^n P(x_j | parents X_j)}
\end{align*}

All terms in the product in the denominator that do not contain xi can be moved outside the summation 
and canceled with the corresponding terms in the numerator. 
The remaining terms in the denominator are those that contain xi, and they can be expressed as a function of the Markov blanket of xi. Thus, we have:
\begin{align*}
P(x_i | x_1, \ldots, x_{i-1}, x_{i+1}, \ldots, x_n) &= \frac{\Pi_{j=1}^n P(x_j | parents X_j)}{\Sigma_{x_i} \Pi_{j=1}^n P(x_j | parents X_j)} \\
&= \frac{P(x_i | parents X_i) \Pi_{Y_j \in Children(X_i)} P(y_j | parents Y_j)}{\Sigma_{x_i} P(x_i | parents X_i) \Pi_{Y_j \in Children(X_i)} P(y_j | parents Y_j)}
\end{align*}

Compare this with the equation (14.12), we can see that 
$P(X_i | X_1, \ldots, X_{i-1}, X_{i+1}, \ldots, X_n) = P(X_i | MB(X_i))$.
which means that the variable \(X\) is independent of all other variables in the network given its Markov blanket.
\end{solution}
%------------------------------%


\newpage










%------------------------------%
\end{document}

